% =============================================================================
% CAPÍTULO 7 — CONCLUSIONES GLOBALES Y MÉTRICAS FINALES
% =============================================================================
\chapter{Conclusiones Globales y Métricas Finales}
\label{ch:conclusiones}
\resumenCapitulo{Consolidación de los resultados de verificación y validación de los cinco
sprints: calidad del producto final de la Release 2.0.0, evaluación de la deuda técnica
acumulada y cierre formal con las firmas de aprobación del equipo de aseguramiento de
calidad.}

% =============================================================================
\section{Resumen de Calidad del Producto Final (Versión Release)}
\label{sec:concl-resumen}

A lo largo de las cinco iteraciones se diseñaron y ejecutaron \textbf{200 casos de prueba}
(40 por sprint) y se registraron \textbf{48 defectos} documentados con calidad forense. La
totalidad de los defectos fue corregida y verificada dentro de su iteración, lo que arroja
una \textbf{eficiencia de remoción de defectos del 100\%} y \textbf{cero defectos críticos
abiertos} al cierre de la versión.

\subsection{Consolidado de Ejecución por Sprint}

\vspace{0.1cm}
\begin{center}
\renewcommand{\arraystretch}{1.35}
\fontsize{8.5}{10.2}\selectfont\sffamily
\begin{tabular}{|c|c|c|c|c|c|c|}
\hline
\rowcolor{headerTabla}
\sffamily\textbf{Sprint} & \sffamily\textbf{Plan.} & \sffamily\textbf{Ejec.} & \sffamily\textbf{PASS} & \sffamily\textbf{FAIL} & \sffamily\textbf{Aprob.} & \sffamily\textbf{Defectos} \\
\hline
Sprint 1 & 40 & 40 & 33 & 6 & 82,5\% & 9 \\ \hline
\rowcolor{grisTecnico}
Sprint 2 & 40 & 40 & 32 & 8 & 80,0\% & 9 \\ \hline
Sprint 3 & 40 & 40 & 29 & 11 & 72,5\% & 11 \\ \hline
\rowcolor{grisTecnico}
Sprint 4 & 40 & 40 & 30 & 10 & 75,0\% & 10 \\ \hline
Sprint 5 & 40 & 40 & 31 & 9 & 77,5\% & 9 \\ \hline
\rowcolor{headerTabla}
\sffamily\textbf{Total} & \sffamily\textbf{200} & \sffamily\textbf{200} & \sffamily\textbf{155} & \sffamily\textbf{44} & \sffamily\textbf{77,5\%} & \sffamily\textbf{48} \\ \hline
\end{tabular}
\end{center}

\begin{NotaTecnica}
La columna \textbf{Aprob.} es la tasa de aprobación en la \textit{primera} ejecución. Tras
el retesteo de los defectos corregidos, los 44 casos en estado FAIL pasaron a PASS, de modo
que la \textbf{tasa de aprobación efectiva al cierre de cada sprint fue del 100\%} (un caso,
TC-S1-039, estuvo temporalmente bloqueado y se verificó tras corregir su dependencia).
\end{NotaTecnica}

\subsection{Tendencia de la Tasa de Aprobación}

La figura muestra la evolución de la tasa de aprobación en primera ejecución. El valle del
Sprint 3 corresponde a la mayor superficie de integración externa y seguridad, recuperado
en los sprints posteriores gracias al refuerzo OWASP del Sprint 4.

\begin{center}
\begin{tikzpicture}[font=\sffamily\scriptsize]
  \def\h{4.0}
  \draw[grisISO,->] (0,0) -- (11,0) node[right]{Sprint};
  \draw[grisISO,->] (0,0) -- (0,\h) node[above]{\% aprobación};
  \foreach \y/\lbl in {0/0,1.25/50,2.0/70,2.5/80,3.0/90}{\draw[grisISO!40] (0,\y*\h/3.5) -- (10.5,\y*\h/3.5); \node[left,text=grisISO,font=\sffamily\tiny] at (0,\y*\h/3.5) {\lbl};}
  % puntos: 82.5, 80, 72.5, 75, 77.5  (escala: %/90*h aprox -> usar y=val/100*h/0.95)
  \def\sx{2.0}
  \coordinate (p1) at (1,{82.5/100*\h});
  \coordinate (p2) at (3,{80.0/100*\h});
  \coordinate (p3) at (5,{72.5/100*\h});
  \coordinate (p4) at (7,{75.0/100*\h});
  \coordinate (p5) at (9,{77.5/100*\h});
  \draw[colorPrimario,very thick] (p1)--(p2)--(p3)--(p4)--(p5);
  \foreach \p/\lbl/\v in {p1/S1/82{,}5,p2/S2/80{,}0,p3/S3/72{,}5,p4/S4/75{,}0,p5/S5/77{,}5}{
    \fill[colorAcento] (\p) circle (2.2pt);
    \node[below=5pt,text=grisISO,font=\sffamily\tiny] at (\p) {\lbl};
    \node[above=3pt,text=colorPrimario,font=\sffamily\tiny\bfseries] at (\p) {\v\%};}
\end{tikzpicture}
\end{center}

\clearpage
\subsection{Distribución Global de Defectos por Severidad}

\vspace{0.1cm}
\noindent
\renewcommand{\arraystretch}{1.35}
\fontsize{8.5}{10.2}\selectfont\sffamily
\begin{tabularx}{\textwidth}{|Z|Z|Z|Z|Z|Z|Z|Z|}
\hline
\rowcolor{headerTabla}
\sffamily\textbf{Severidad} & \sffamily S1 & \sffamily S2 & \sffamily S3 & \sffamily S4 & \sffamily S5 & \sffamily\textbf{Total} & \sffamily\textbf{\% del total} \\
\hline
\sevCrit & 2 & 1 & 3 & 3 & 3 & 12 & 25,0\% \\ \hline
\rowcolor{grisTecnico}
\sevAlt & 3 & 3 & 4 & 4 & 2 & 16 & 33,3\% \\ \hline
\sevMed & 3 & 4 & 4 & 3 & 3 & 17 & 35,4\% \\ \hline
\rowcolor{grisTecnico}
\sevBaj & 1 & 1 & 0 & 0 & 1 & 3 & 6,3\% \\ \hline
\rowcolor{headerTabla}
\sffamily\textbf{Total} & 9 & 9 & 11 & 10 & 9 & \sffamily\textbf{48} & \sffamily\textbf{100\%} \\ \hline
\end{tabularx}

\vspace{0.3cm}
\noindent La concentración de defectos críticos en los sprints 3, 4 y 5 es coherente con la
incorporación de la superficie de integración externa, los controles OWASP y el panel de
administración; en todos los casos fueron de seguridad (control de acceso, SSRF, XSS,
revocación de tokens y exposición de secretos) y se resolvieron antes del cierre.

\subsection{Defectos por Componente (Mapa de Calor)}

\vspace{0.1cm}
\noindent
\renewcommand{\arraystretch}{1.3}
\fontsize{8.5}{10.2}\selectfont\sffamily
\begin{tabularx}{\textwidth}{|>{\ttfamily\bfseries\color{colorPrimario}}l|c|Y|}
\hline
\rowcolor{headerTabla}
\sffamily\textbf{Componente} & \sffamily\textbf{Defectos} & \sffamily\textbf{Observación} \\
\hline
security & 8 & JWT, rate limiting, blacklist, CSRF: foco principal de hallazgos críticos. \\ \hline
\rowcolor{grisTecnico}
frontend & 11 & Accesibilidad, estado, i18n, tema y persistencia de UI. \\ \hline
auth & 5 & Registro, OTP, OAuth sync y enumeración de usuarios. \\ \hline
\rowcolor{grisTecnico}
connections & 3 & SSRF y ciclo de vida de tokens de proveedor. \\ \hline
chat / cvstudio & 5 & Aislamiento de canales, idempotencia, timeouts e IA. \\ \hline
\rowcolor{grisTecnico}
portfolio & 4 & Visibilidad, slug y protección con contraseña. \\ \hline
skills / experience & 6 & IDOR, integridad referencial, validaciones de dominio. \\ \hline
\rowcolor{grisTecnico}
admin / metrics & 4 & Control de acceso, logs y rendimiento del dashboard. \\ \hline
db / uploads & 2 & Seeders, índices y validación de contenido. \\ \hline
\end{tabularx}

\subsection{Cobertura de Pruebas Consolidada}

\barraCobertura{82}{Cobertura media de líneas del backend (Java / jOOQ)}
\barraCobertura{77}{Cobertura media de ramas del backend (validaciones y seguridad)}
\barraCobertura{70}{Cobertura media de componentes del frontend (Vitest + Testing Library)}

\vspace{0.2cm}
\noindent La cobertura se mantuvo por encima de los umbrales mínimos acordados (80\% de
líneas en backend y 65\% de componentes en frontend) en todos los sprints, complementada por
la verificación estática (ESLint y TypeScript estricto) como primera línea de defensa.

\subsection{Lecciones Aprendidas del Proceso de V\&V}

El análisis transversal de los 48 defectos arroja patrones útiles para el mantenimiento y
para futuros proyectos del equipo:

\begin{itemize}
    \item \textbf{La seguridad concentra el riesgo real.} El 25\% de los defectos fueron
    críticos y todos de seguridad (control de acceso, SSRF, XSS, tokens). Adelantar las
    pruebas de seguridad a la fase de diseño habría reducido su coste de corrección.
    \item \textbf{La superficie externa es la más frágil.} Las integraciones (conexiones,
    subida de archivos, IA, tiempo real) generaron la mayor densidad de defectos del
    proyecto (Sprint 3), confirmando el valor de las pruebas basadas en riesgo.
    \item \textbf{Las pruebas con dependencias reales detectan más.} Testcontainers reveló
    defectos de integridad (cascadas, índices, N+1) que una base de datos en memoria habría
    ocultado.
    \item \textbf{El control de acceso debe probarse por rol y por propietario.} Los IDOR y
    la fuga de métricas se detectaron solo al probar explícitamente accesos cruzados.
    \item \textbf{La calidad transversal (a11y, i18n) requiere verificación automatizada.}
    Los defectos de accesibilidad e internacionalización motivaron incorporar comprobaciones
    en CI (deuda DT-02 y DT-04).
\end{itemize}

% =============================================================================
\clearpage
\section{Evaluación de la Deuda Técnica Acumulada}
\label{sec:concl-deuda}

Aunque no quedaron defectos críticos ni altos abiertos, la V\&V identificó elementos de
\textbf{deuda técnica} y mejoras diferidas que, sin bloquear la \textit{Release} 2.0.0, se
recomienda atender en el mantenimiento evolutivo. Ninguno representa un riesgo de seguridad
residual.

\vspace{0.1cm}
\noindent
\renewcommand{\arraystretch}{1.3}
\fontsize{8.5}{10.2}\selectfont\sffamily
\begin{tabularx}{\textwidth}{|>{\ttfamily\bfseries\color{colorPrimario}}p{1.8cm}|Y|c|}
\hline
\rowcolor{headerTabla}
\sffamily\textbf{ID} & \sffamily\textbf{Descripción de la deuda / mejora diferida} & \sffamily\textbf{Prioridad} \\
\hline
DT-01 & Ampliar la cobertura E2E automatizada (hoy parcialmente manual) con una suite de extremo a extremo en CI. & \sevMed \\ \hline
\rowcolor{grisTecnico}
DT-02 & Externalizar la verificación de cobertura i18n a una validación de CI que falle ante claves faltantes. & \sevMed \\ \hline
DT-03 & Incorporar pruebas de carga sostenida (más allá del rendimiento puntual) sobre chat y dashboards. & \sevMed \\ \hline
\rowcolor{grisTecnico}
DT-04 & Auditoría de accesibilidad por un tercero independiente para certificación WCAG AA formal. & \sevBaj \\ \hline
DT-05 & Monitorización continua de SSRF/egress con lista de permitidos centralizada y observabilidad. & \sevMed \\ \hline
\rowcolor{grisTecnico}
DT-06 & Telemetría de errores en producción (correlación de \textit{stacktraces}) para acelerar el RCA. & \sevBaj \\ \hline
DT-07 & Migrar las pruebas de seguridad ad-hoc (SQLi, XSS, JWT) a una suite DAST recurrente. & \sevMed \\ \hline
\end{tabularx}

\vspace{0.3cm}
\noindent\textbf{Riesgos residuales aceptados.} Los servicios externos (Supabase, Gemini,
Google Drive) quedan gobernados por sus respectivos SLA; su disponibilidad y seguridad
internas están fuera del control directo de Byte Busters S.R.L. y se mitigan con
degradación elegante, \textit{timeouts} y validación de egress.

\begin{AvisoNota}
La deuda técnica listada se incorporó al \textit{Product Backlog} de mantenimiento. Ningún
elemento condiciona la aceptación de la \textit{Release} 2.0.0: todos son mejoras de
robustez, observabilidad o automatización del proceso de calidad.
\end{AvisoNota}

% =============================================================================
\clearpage
\section{Cierre y Firmas de Aprobación}
\label{sec:concl-cierre}

Sobre la base de la evidencia consolidada —200 casos ejecutados, 48 defectos remediados al
100\%, cero defectos críticos abiertos, cobertura por encima de umbrales y regresión final
en verde— el equipo de aseguramiento de calidad de Byte Busters S.R.L. emite el siguiente
veredicto global.

\begin{Veredicto}
\textbf{PRODUCTO APROBADO PARA LA RELEASE 2.0.0.} La plataforma EthosHub satisface los
requisitos funcionales y no funcionales verificados y validados a lo largo de los cinco
sprints. Se certifica su aptitud para el despliegue en producción, con la deuda técnica
registrada para el mantenimiento evolutivo.
\end{Veredicto}

\vspace{0.4cm}
\noindent
\renewcommand{\arraystretch}{2.2}
\fontsize{9}{11}\selectfont\sffamily
\begin{tabularx}{\textwidth}{|>{\bfseries\color{colorPrimario}}p{4.2cm}|Y|Y|}
\hline
\rowcolor{headerTabla}
\sffamily\textbf{Rol} & \sffamily\textbf{Nombre} & \sffamily\textbf{Firma / Fecha} \\
\hline
QA Lead & C.B. Quispe Flores & \hrulefill \\ \hline
\rowcolor{grisTecnico}
Scrum Master / Arquitecto & J.D. Pardo Pozo & \hrulefill \\ \hline
Representante Legal & Juan Diego Pardo Pozo & \hrulefill \\ \hline
\rowcolor{grisTecnico}
Consultor TIS supervisor & Corina J. Flores Villarroel & \hrulefill \\ \hline
\end{tabularx}

\vspace{0.6cm}
\begin{center}
{\sffamily\footnotesize\color{grisISO}
Documento \textbf{QA-VV-ETHOSHUB-2.0.0} \,$\cdot$\, Estado: Aprobado \,$\cdot$\,
Cochabamba, Bolivia \,$\cdot$\, Junio de 2026\par}
\end{center}
